\chapter{Work Remaining}

\section{Model Evaluation}
Rigorous testing is required to assess the effectiveness of our recommendation system. Evaluation will focus on accuracy metrics, user satisfaction, scalability testing.
Evaluation results will drive improvements to enhance recommendation accuracy and user engagement.

\section{User Interface (UI) Enhancement}
A seamless and intuitive user interface is crucial for maximizing user experience. The following UI improvements are planned:

\begin{itemize}
    \item Develop an interactive UI to display recommended attractions with essential details (descriptions, images, user ratings).
    \item Implement filtering, sorting, and search options to allow users to customize recommendations based on preferences such as budget, location, or interest.
    \item Ensure responsiveness across devices (desktop, tablet, smartphone) for seamless access.
\end{itemize}

These changes will enhance the accessibility and usability of the recommendation system, improving user engagement.

\section{Map Integration and Optimization}
The map feature is an essential component of the recommendation system, and improvements are needed:

\begin{itemize}
    \item \textbf{Interactive Map Features}: Enable users to select attractions from the map and dynamically adjust preferences (e.g., filtering by type, distance.).
    %\item \textbf{Real-Time Data Integration}: Incorporate real-time data (e.g., transportation options, traffic conditions) for accurate travel suggestions.
\end{itemize}

These enhancements will allow users to plan more efficient and enjoyable trips.

\section{Deployment and Maintenance}
The final phase involves deploying the recommendation system and ensuring its ongoing functionality:

\begin{itemize}
    \item \textbf{Deployment Infrastructure}: Set up servers and databases for hosting the system, ensuring scalability and handling user traffic.
    \item \textbf{User Feedback Mechanism}: Establish a system to collect user feedback, enabling continuous improvement.
\end{itemize}

\section{Future Improvements}
Looking ahead, several areas for improvement are identified:

\begin{itemize}
    \item \textbf{Personalized Notifications}: Push notifications to suggest new attractions or events based on user preferences and location.
    \item \textbf{Advanced Collaborative Filtering}: Incorporate deep learning techniques to enhance personalization in the collaborative filtering model.
    \item \textbf{Social Media Integration}: Enable users to share recommendations and experiences on social media, enhancing engagement and providing additional data for the system.
\end{itemize}

These future improvements will further personalize the experience, making the recommendation system more dynamic and interactive.

By systematically addressing these areas, we aim to deliver a comprehensive, user-centric recommendation system that adapts to user preferences, enhances the travel experience, and evolves over time based on feedback and new data.
